Modern science delegates much of its quality certification to a small set of prestigious venues. We ask how rapidly the signaling value of acceptance decays as submission volume grows, and through what mechanisms. We formalize gate-keeping as a capacity-constrained Bayesian signaling problem in which $\Nsub$ papers compete for $\Kcap$ slots under reviewer noise $\sigrev$, and we instantiate the model at three layers: paper (mutual information $\miquant$ and tail recall), reader (consumption under a fixed time budget), and researcher (publication count as a proxy for ability). On the analytical side we show that under fixed capacity, mutual information and recall on the true top-1\% of papers fall monotonically in $\Nsub$ for every reviewer noise level (Kendall $\tau\le-0.95$, $p<10^{-6}$), and that the Spearman correlation between researcher ability and publication count drops from $0.78$ at $M{=}200$ researchers to $0.26$ at $M{=}20{,}000$. Under linearly expanding capacity ($\Kcap = \arate \cdot \Nsub$) the per-paper signal is preserved but the reader's budget cannot keep up, and load-induced reviewer noise re-introduces the same decay. We confirm the prediction empirically on $14$ years of \iclr submissions: the standardized posterior contrast $\hcontrast/\stdrating$ falls log-linearly in $\Nsub$ with slope $-0.32$ per $e$-fold ($R^2=0.85$, $p=1.1\times10^{-3}$), from $2.55$ at $\Nsub{=}490$ to $1.08$ at $\Nsub{=}19{,}814$, while mean per-paper reviewer SD grows from $0.86$ to $1.53$. The model recovers Akerlof, Spence, and Heckman--Moktan as limiting cases and yields a sharp policy implication: no capacity policy simultaneously preserves per-paper precision, researcher identifiability, and reader-level relevance.
